\documentclass[10pt,a4paper]{amsart}
\usepackage[margin=19mm]{geometry}
\usepackage{amsmath,amssymb,mathtools}
\usepackage[scr=rsfs,cal=euler]{mathalfa}
\usepackage{xfrac}
\usepackage[unicode,hidelinks,bookmarks=false]{hyperref}
\newcommand{\ie}{i.e.\ }
\newcommand{\cf}{cf.\ }
\newcommand{\resp}{respectively\ }
\newcommand{\Subsection}{\subsection}
\providecommand{\nobreakdash}{\nobreak\hbox{-}}
\setlength{\emergencystretch}{2em}
\multlinegap0pt
\usepackage{bm}
\usepackage{amssymb}
\usepackage{dsfont}
\usepackage{float}
\usepackage{xcolor}
\newtheorem{proposition}{Proposition}
\newtheorem{theorem}{Theorem}
\title{\bf Conditional copula representations and extremal bounds for multivariate statistical functionals}
\author{Roberto Vila, Cira E G Otiniano, Carolyne Brito, Enzo Brasil}
\date{Source: arXiv:2607.26256v1 (CC BY 4.0)}
\begin{document}
\maketitle

\noindent\emph{Source and licence.} \bf Conditional copula representations and extremal bounds for multivariate statistical functionals. Version arXiv:2607.26256v1 (CC BY 4.0), distributed by arXiv under the Creative Commons Attribution 4.0 International licence, http://creativecommons.org/licenses/by/4.0/, as recorded in the arXiv licence metadata for this exact version and shown on the abstract page https://arxiv.org/abs/2607.26256v1. Full licence notice: \texttt{licenses/arxiv-2607.26256-CC-BY-4.0.txt}.\par\smallskip

\noindent\textit{Editorial scope.} Counted unit: the article's theorem main-result-multivariate (source0.tex lines 1045--1085). The article's complete original proof of it is delivered with it (source0.tex lines 1087--1131, one \texttt{proof} environment of 45 lines). No mathematics is written, repaired or re-derived: the statement and the proof are the article's own lines, in the article's own notation, and no retained line is substituted or re-flowed.  Retained uncounted as the source-local context the proof needs: lines 854--874; lines 994--1037.  Nothing was omitted from any retained span, so the package carries no omission record.  No retained line contains a citation, so the wrapper carries no bibliography and no bibliography entry.  No retained line contains a figure environment, an image-inclusion command or drawing code, so no image is delivered and the package declares no asset dependency.  Editorial formatting only, in addition: an AMS \texttt{amsart} preamble that restates the article's own declarations and macros used by the retained lines, namely the environments \texttt{proposition}, \texttt{theorem} on separate counters, together with the standard packages those lines need.  The retained environments print in this document as Proposition 1, Theorem 1 in order of appearance, whereas the article prints the same items with the numbering of its own full document.  The complete native source of the article as distributed in the arXiv e-print for this exact version is bundled unchanged as \texttt{sources/206249/source0.tex}.
\begin{align}
	\pi_g(C)
	\equiv 
	\mathbb E_C[g(\boldsymbol X)]
	&=
	\int_{\mathbb R^n}
	g(x_1,\ldots,x_n)\,
	{\rm d}
	C\!\left(
	F_1(x_1),\ldots,F_n(x_n)
	\right)
	\nonumber
	\\[0.2cm]
	&=
	\int_{[0,1]^n}
	g\!\left(
	F_1^{-1}(u_1),\ldots,F_n^{-1}(u_n)
	\right)
	{\rm d}C(u_1,\ldots,u_n),
	\label{id-exp-cop}
\end{align}

\begin{proposition}\label{antitonic}
	Let
	$g:\mathbb R^n\to\mathbb R$
	be
	$\Delta$-antitonic.
	Then the functional
	$\pi_g$
	is increasing with respect to the concordance order. Hence, for every
	$C\in\mathcal C_n$,
	\[
	\inf_{D\in\mathcal C_n}\pi_g(D)
	\leqslant
	\pi_g(C)
	\leqslant
	\pi_g(M_n),
	\]
	where $\mathcal C_n$ denotes the class of all $n$-copulas 
	
	Moreover, if
	$-g$
	is
	$\Delta$-antitonic,
	then
	\[
	\pi_g(M_n)
	\leqslant
	\pi_g(C)
	\leqslant
	\sup_{D\in\mathcal C_n}\pi_g(D).
	\]
	When $n=2$, the bounds are attained and
	$
	\pi_g(W_2)
	=
	\inf_{D\in\mathcal C_2}\pi_g(D)
	$
	in the first case, whereas
	$
	\pi_g(W_2)
	=
	\sup_{D\in\mathcal C_2}\pi_g(D)
	$
	in the second case.
\end{proposition}

\begin{theorem}\label{main-result-multivariate}
	Let
	$g:\mathbb R^n\to\mathbb R$
	be
	$\Delta$-antitonic.
	Then
	\[
	\inf_{D\in\mathcal C_n}\pi_g(D)
	\leqslant
	\mathbb E[g(\boldsymbol X)]
	\leqslant
	\int_0^1
	g\!\left(
	F_1^{-1}(u),
	\ldots,
	F_n^{-1}(u)
	\right)
	{\rm d}u.
	\]
	
	Moreover, if
	$-g$
	is
	$\Delta$-antitonic,
	then
	\[
	\int_0^1
	g\!\left(
	F_1^{-1}(u),
	\ldots,
	F_n^{-1}(u)
	\right)
	{\rm d}u
	\leqslant
	\mathbb E[g(\boldsymbol X)]
	\leqslant
	\sup_{D\in\mathcal C_n}\pi_g(D).
	\]
	
	In particular, when $n=2$, the lower (respectively, upper) bound is attained at the Fréchet--Hoeffding lower copula $W_2$.
\end{theorem}

\begin{proof}
	Since
	$
	\mathbb E[g(\boldsymbol X)]
	=
	\pi_g(C),
	$
	Proposition~\ref{antitonic} yields
	$$
	\mathbb E[g(\boldsymbol X)]
	\leqslant
	\pi_g(M_n).
	$$
	By \eqref{id-exp-cop},
	\[
	\pi_g(M_n)
	=
	\int_{[0,1]^n}
	g\!\left(
	F_1^{-1}(u_1),
	\ldots,
	F_n^{-1}(u_n)
	\right)
	{\rm d}M_n(u_1,\ldots,u_n).
	\]
	Since
	$
	M_n(u_1,\ldots,u_n)
	=
	\min\{u_1,\ldots,u_n\},
	$
	the probability measure induced by
	$M_n$
	is concentrated on the diagonal
	$
	\{(u,\ldots,u):0<u<1\}.
	$
Equivalently, \[ \frac{\partial^n M_n(u_1,\ldots,u_n)} {\partial u_1\cdots\partial u_n} = \delta_u(u_2)\cdots\delta_u(u_n) \] in the sense of distributions. Therefore, 
\[ 
\pi_g(M_n) = \int_0^1 g\!\left( F_1^{-1}(u), \ldots, F_n^{-1}(u) \right) {\rm d}u, 
\] 
which proves the first inequality.
	
	The second inequality follows by applying the first part to the function $-g$.
\end{proof}

\end{document}
